\section{La potencia como producto repetido}

Considere nuevamente el producto
\[
2\cdot2\cdot2\cdot2.
\]
Como los cuatro factores son iguales, podemos abreviarlo mediante
\[
2^4.
\]
Por lo tanto,
\[
2^4=2\cdot2\cdot2\cdot2=16.
\]
Las tres escrituras representan el mismo número, pero muestran aspectos diferentes. La primera hace visibles los factores; la segunda señala directamente qué factor se repite y cuántas veces; la tercera es el numeral obtenido al evaluar la expresión.

\begin{definitionbox}[Potencia con exponente natural positivo]
Sean \(a\in\mathbb N\) y \(n\in\mathbb N\) con \(n>0\). Se define
\[
a^n=
\underbrace{a\cdot a\cdot\ldots\cdot a}_{n\text{ factores}}.
\]
El número \(a\) recibe el nombre de \emph{base}; el número \(n\), el de \emph{exponente}; y el valor de la expresión completa, el de \emph{potencia}.
\end{definitionbox}

Por ejemplo, en
\[
3^5=3\cdot3\cdot3\cdot3\cdot3=243,
\]
la base es \(3\), el exponente es \(5\) y la potencia es \(243\).

La posición de cada número es parte de la notación. La base indica cuál es el factor repetido; el exponente indica cuántas veces aparece. Por eso, en general, intercambiarlos modifica la cantidad representada:
\[
2^3=2\cdot2\cdot2=8,
\qquad
3^2=3\cdot3=9.
\]
No existe una propiedad conmutativa de la potenciación: normalmente \(a^n\ne n^a\).

\subsection{Exponente uno}

Si el exponente es \(1\), la base aparece como factor una sola vez. En consecuencia,
\[
a^1=a.
\]
No es necesario imaginar un producto de varios factores: la expresión \(a^1\) representa simplemente el mismo número que \(a\).

\subsection{Exponente cero}

La definición como producto repetido nos permite entender las potencias con exponente positivo. También muestra una relación sencilla entre potencias consecutivas.

Por ejemplo,
\[
2^4
=2\cdot2\cdot2\cdot2
=(2\cdot2\cdot2)\cdot2
=2^3\cdot2.
\]
Del mismo modo,
\[
2^3=2^2\cdot2,
\qquad
2^2=2^1\cdot2.
\]

En general, aumentar el exponente en una unidad equivale a agregar un factor igual a la base:
\[
a^{n+1}=a^n\cdot a,
\qquad n>0.
\]
Esta igualdad no es una nueva regla que debamos memorizar; se obtiene directamente de la definición de potencia como producto repetido.

Ahora queremos extender la notación para incluir el exponente \(0\). Para hacerlo de manera compatible con la relación anterior, debería cumplirse también
\[
a^1=a^0\cdot a.
\]
Pero ya sabemos que
\[
a^1=a.
\]
Por lo tanto,
\[
a=a^0\cdot a.
\]

Si \(a\ne0\), podemos dividir ambos lados entre \(a\):
\[
\frac{a}{a}
=
\frac{a^0\cdot a}{a},
\]
y obtenemos
\[
1=a^0.
\]

Esto nos conduce a la siguiente definición.

\begin{definitionbox}[Potencia con exponente cero]
Para todo número \(a\ne0\), se define
\[
a^0=1.
\]
\end{definitionbox}

La condición \(a\ne0\) es importante. Si intentáramos hacer el mismo razonamiento con base cero, la relación
\[
a^1=a^0\cdot a
\]
se convertiría en
\[
0=0^0\cdot0.
\]
Pero esta igualdad se cumple cualquiera sea el valor colocado en lugar de \(0^0\), ya que todo número multiplicado por cero da cero. Por lo tanto, este razonamiento no determina un valor para \(0^0\).

En este libro dejaremos \(0^0\) sin definir. En contextos más avanzados puede asignársele algún valor por conveniencia, pero esa elección depende del contexto.

En cambio, cuando el exponente es positivo,
\[
0^n=0
\qquad(n>0),
\]
porque el producto contiene al menos un factor igual a cero.

\subsection{La base no tiene que natural}

Hasta ahora hemos utilizado principalmente números positivos como base, pero la idea de potencia como producto repetido también funciona cuando la base es negativa.

Por ejemplo,
\[
(-2)^3=(-2)\cdot(-2)\cdot(-2)=-8,
\]
mientras que
\[
(-2)^4=(-2)\cdot(-2)\cdot(-2)\cdot(-2)=16.
\]

La diferencia está en la cantidad de factores negativos.

Recordemos que el producto de dos números negativos es positivo:
\[
(-a)(-a)=a^2.
\]

Por eso, si tenemos una cantidad par de factores negativos, podemos agruparlos de dos en dos. Cada pareja produce un resultado positivo. Por ejemplo,
\[
(-2)^6
=
\bigl((-2)(-2)\bigr)
\bigl((-2)(-2)\bigr)
\bigl((-2)(-2)\bigr),
\]
y por lo tanto el resultado es positivo.

En cambio, si la cantidad de factores es impar, después de formar todas las parejas posibles queda un factor negativo sin pareja:
\[
(-2)^5
=
\bigl((-2)(-2)\bigr)
\bigl((-2)(-2)\bigr)
(-2),
\]
de modo que el resultado es negativo.

Este razonamiento no depende de cuántos factores haya. Siempre podemos agrupar los factores negativos de dos en dos. Por ello, para \(a>0\),
\[
(-a)^n=
\begin{cases}
a^n, & \text{si \(n\) es par},\\[4pt]
-a^n, & \text{si \(n\) es impar}.
\end{cases}
\]

No es necesario memorizar esta expresión: el signo puede determinarse simplemente observando si la cantidad de factores negativos es par o impar.

\begin{note}
  Mucho cuidado con expresiones como 
  \[
    -a^n
  \]
  ya que el signo \(-\) en tal caso no se considera parte de la potencia. Es decir, la interpretación es 
  \[
    - (a^n)
  \]
\end{note}

\subsubsection{Los paréntesis importan}

Cuando la base es negativa, los paréntesis permiten indicar que el signo forma parte de la base.

Por ejemplo,
\[
(-2)^4
\qquad\text{y}\qquad
-2^4
\]
representan expresiones diferentes.

En
\[
(-2)^4,
\]
la base es \(-2\), por lo que
\[
(-2)^4
=
(-2)(-2)(-2)(-2)
=
16.
\]

En cambio, en
\[
-2^4,
\]
la base de la potencia es \(2\). Primero se calcula
\[
2^4=16
\]
y luego se toma su opuesto:
\[
-2^4=-(2^4)=-16.
\]

Por lo tanto,
\[
(-2)^4=16,
\qquad
-2^4=-16.
\]

\subsection{Una fracción también puede ser la base}

La base de una potencia también puede ser una fracción. La interpretación sigue siendo exactamente la misma: elevar a un exponente natural positivo significa multiplicar la base por sí misma tantas veces como indica el exponente.

Por ejemplo,
\[
\left(\frac{2}{3}\right)^3
=
\frac{2}{3}\cdot
\frac{2}{3}\cdot
\frac{2}{3}.
\]

Como ya sabemos multiplicar fracciones, podemos multiplicar los numeradores entre sí y los denominadores entre sí:
\[
\left(\frac{2}{3}\right)^3
=
\frac{2\cdot2\cdot2}{3\cdot3\cdot3}
=
\frac{8}{27}.
\]

Observemos ahora la forma en que aparecen el numerador y el denominador:
\[
2\cdot2\cdot2=2^3,
\qquad
3\cdot3\cdot3=3^3.
\]
Por lo tanto,
\[
\left(\frac{2}{3}\right)^3
=
\frac{2^3}{3^3}.
\]

El mismo razonamiento funciona con cualquier fracción cuyo denominador sea distinto de cero. Si \(b\ne0\), entonces
\[
\left(\frac{a}{b}\right)^n
=
\frac{a^n}{b^n},
\qquad n>0.
\]

Esta igualdad no introduce una nueva manera de calcular potencias. Es simplemente lo que obtenemos al escribir la potencia como un producto repetido:
\[
\left(\frac{a}{b}\right)^n
=
\underbrace{
\frac{a}{b}\cdot
\frac{a}{b}\cdots
\frac{a}{b}
}_{n\text{ factores}}
=
\frac{
  \overbrace{a\cdot a\cdots a}^{n\text{ factores}}
}{
\underbrace{b\cdot b\cdots b}_{n\text{ factores}}
}
=
\frac{a^n}{b^n}.
\]

La fracción también puede ser negativa. Por ejemplo,
\[
\left(-\frac{2}{3}\right)^3
=
\left(-\frac{2}{3}\right)
\left(-\frac{2}{3}\right)
\left(-\frac{2}{3}\right)
=
-\frac{8}{27},
\]
mientras que
\[
\left(-\frac{2}{3}\right)^4
=
\frac{16}{81}.
\]

Nuevamente, el signo depende de la cantidad de factores negativos: un exponente par produce un resultado positivo y un exponente impar produce un resultado negativo.
